\documentclass[14pt]{extarticle}
\usepackage[T1]{fontenc}
\usepackage[margin=1in]{geometry}
\usepackage{amsthm,amsmath,amssymb}

\newtheorem{exercise}{Exercise}[section]
\setcounter{section}{9}

\newcommand{\inv}[1]{#1^{-1}}
\newcommand{\join}[3][,]{#2_0 #1 #2_1 #1 \cdots #1 #2_{#3}}
\newcommand{\N}{\mathbb{N}}
\DeclareMathOperator{\Abelian}{Abelian}

\begin{document}

\setcounter{exercise}{10}
\begin{exercise}
  Prove that a quotient group of a cyclic group is cyclic.
\end{exercise}
\begin{proof}
  For any cyclic group $G$ and normal subgroup $H$,
  let $G = \langle g \rangle$ and $H = \langle g^n \rangle$ for some minimum $n \in \N$.
  We claim $G/H \approx Z_n$ or $G/H \approx G$ if $n = 0$ which is trivial.

  We claim:
  {
    \newcommand{\gH}{\langle gH \rangle}
    \newcommand{\GH}{G/H}

    \begin{center}
      \boxed{\gH = \GH}
    \end{center}


    Every element in $\gH$ is a coset of $H$, thus $\gH \subseteq \GH$.

    Any element in $\GH$ has form $hH$ where $h \in G$, therefore $h = g^s$ for some $s$.
    Then $hH = g^sH \in \langle gH \rangle$.

    We claim $| gH | = n$.
    $(gH)^n = g^nH$ where $g^n \in H$, so $g^nH = H$.
    Suppose $0 < m < n, g^mH = H$. Then $g^m \in H$ and $g^{\gcd(m, n)} \in H$
    where $\gcd(m, n) \le m$ and $\gcd(m, n)$ divides $n$.
    But this contradict our assumption that $n$ is minimum such that $\langle g^n \rangle = H$,
    because $\langle g^{\gcd(m, n)} \rangle = H$.

    Thus, $\gH$ is cylic and order $n$, which is isomorphic to $Z_n$
  }
\end{proof}

\begin{exercise}
  Prove that a quotient group of an Abelian group is abelian.
\end{exercise}
\begin{proof}
  For any Abelian group $G$ and normal subgroup $H$. For all $aH$ and $bH$ in $G/H$
  where $a \ b \in G$.
  $(aH)(bH) = abH = baH = (bH)(aH)$.
\end{proof}

\setcounter{exercise}{20}

\begin{exercise}
  For any Abelian group $G$ of order $\join[]{p}{n - 1}$ where $p_i$ are distinct primes.
  Shows that $G$ is cyclic.
\end{exercise}
\begin{proof}
  Since $|G| = \join[]{p}{n - 1}$, there are elements of each prime orders, say,
  $|g_0| = p_0, |g_1| = p_1 \cdots$.
  We claim $G = \langle g_0 \rangle \times \langle g_1 \rangle \times \cdots \times \langle g_{n - 1} \rangle$.
  
  They are all normal because $G$ is Abelian, so the first property satisfied.
  Let $H = (\langle g_0 \rangle \langle g_1 \rangle \cdots \langle g_{i - 1} \rangle) \cap \langle g_i \rangle$
  for some $i$.
  $H$ must be the subgroup of both $\langle g_0 \rangle \langle g_1 \rangle \cdots \langle g_{i - 1} \rangle$ and $\langle g_i \rangle$,
  therefore, $|H|$ divides $\join[]{p}{i - 1}$ and $p_i$. But all $p$'s are distinct prime, so $|H|$ must be $1$,
  thus $H = \{ e \}$.

  So $G$ is the internal direct product of $\langle g_0 \rangle \times \langle g_1 \rangle \times \cdots \times \langle g_{n - 1} \rangle$,
  which is isomorphic to $G^\prime = \langle g_0 \rangle \oplus \langle g_1 \rangle \oplus \cdots \oplus \langle g_{n - 1} \rangle$.
  And the order of $\langle g_i \rangle$ are relative primes, so $G^\prime$ is cyclic, so is $G$.
\end{proof}

\setcounter{exercise}{40}

\begin{exercise}
  Let $H$ be proper subgroup of $Q$, the group of rational numbers under addition.
  Show that $H$ is infinite index.
\end{exercise}
\begin{proof}
  Since $Q$ Abelian, we need to show $Q/H$ is infinite.
  Suppose $|Q/H|$ is some finite $n$, let $aH \in Q/H$ and $aH \neq H$,
  then $(aH)^n = (na)H = H$.
  But we found that $(\frac{a}{n})H \in Q/H$ since it is a coset of $H$,
  but $((\frac{a}{n})H)^n = aH$ which is not identity, contradict the fact that
  $\forall aH \in G/H , (aH)^n = H$
\end{proof}

\end{document}